% GENERATED FILE. Edit canonical lessons, then run npm run generate:books.
% canonical-source-hash: fnv1a64:3277224f8befd42e
% canonical-lessons: PA-C45-labels, PA-C45-form, PA-C45-pehla-paath

\chapter{Reading — Labels, a Form, and the Same Facts as Sentences}
\label{ch:pa-pehla-paath}
\hlchaptermodality{45}

\begin{chapteropening}
\textbf{By the end of this chapter:} \emph{I can read a Punjabi form label cold, read a form somebody else filled in, and follow a short paragraph without translating word by word.}

The six facts of a form said again as sentences, so the colon the form leans on is replaced by the \pa{ਹੈ} Punjabi puts last.
\end{chapteropening}

\section[nāṁ · umar · kamm · rihāiś · bhāśā …]{(reading) — six labels, no blank after them}

\label{lesson:PA-C45-labels}



\begin{quote}
You have written every one of these into a form. Not one is new.

What is new is that there is no blank after it, and nothing telling you which is coming.
\end{quote}

\begin{tcolorbox}[breakable,skin=enhanced,colback=orange!6,colframe=orange!45!black,boxrule=0.5pt,arc=1mm,left=6pt,right=6pt,top=4pt,bottom=4pt,fonttitle=\bfseries,title={Read this}]
\begin{quote}
\pa{ਨਾਂ} \pa{ਉਮਰ} \pa{ਕੰਮ} \pa{ਰਿਹਾਇਸ਼} \pa{ਭਾਸ਼ਾ} \pa{ਫ਼ੋਨ}
\end{quote}

Read down the list once. Do not assemble each one letter by letter — look at the whole word and let it arrive.

Again, and notice which came at once and which you had to build. Both are fine. The difference is what practice removes.
\end{tcolorbox}

\begin{grammarlens}[title={What you've built}]
Filling a form, the label is the part you can ignore. It sits there, you know what goes next to it, and your attention is on the blank.

Reading one, the label is the whole message. \textbf{\pa{ਰਿਹਾਇਸ਼}} on its own tells you what a stranger is about to be asked.

That is why six is enough for one sitting: the same six signs, doing the opposite job.
\end{grammarlens}

\subsection*{Your turn}
\begin{itemize}
\raggedright
  \item \emph{Say it:} the six labels down the list, once, without stopping
\end{itemize}

\emph{Twice through:}

\begin{tcolorbox}[breakable,colback=teal!4,colframe=teal!35!black,title={Before you move on}]
How many of the six were new? (\textbf{None.})

What was new? (\textbf{There was no blank}, and nobody said them first.)
\end{tcolorbox}
\section[nāṁ: aman · umar: 25 · kamm: naukarī]{(reading) — a form somebody else filled in}

\label{lesson:PA-C45-form}



\begin{quote}
The last lesson gave you the labels alone. Here each one has an answer beside it.

Every value below is one you have written yourself.
\end{quote}

\begin{tcolorbox}[breakable,skin=enhanced,colback=orange!6,colframe=orange!45!black,boxrule=0.5pt,arc=1mm,left=6pt,right=6pt,top=4pt,bottom=4pt,fonttitle=\bfseries,title={Read this}]
\begin{quote}
\pa{ਨਾਂ}: \pa{ਅਮਨ} \pa{ਉਮਰ}: \pa{੨੫} \pa{ਕੰਮ}: \pa{ਨੌਕਰੀ} \pa{ਰਿਹਾਇਸ਼}: \pa{ਸ਼ਹਿਰ} \pa{ਭਾਸ਼ਾ}: \pa{ਪੰਜਾਬੀ} \pa{ਫ਼ੋਨ}: \pa{੦੧੨} \pa{੨੫੧}
\end{quote}

Read it through once, without stopping.

Now answer without looking back: what does this person do?
\end{tcolorbox}

\begin{grammarlens}[title={What you've built}]
A form is the first real thing a beginner can read end to end, because it is built to be read fast and gives you the label before the value every time.

Notice the colon does the work a verb would do in a sentence. There is no \textbf{\pa{ਹੈ}} anywhere on this page, and nothing is missing.

Six lines, and you now know six things about somebody you have never met.
\end{grammarlens}

\subsection*{Your turn}
\begin{itemize}
\raggedright
  \item \emph{Say it:} the six lines, once, in order, without stopping
\end{itemize}

\emph{Twice through:}

\begin{tcolorbox}[breakable,colback=teal!4,colframe=teal!35!black,title={Before you move on}]
Whose form is this? (\textbf{\pa{ਅਮਨ}.})

Which word does the colon replace? (\textbf{\pa{ਹੈ}} — and nothing is missing without it.)
\end{tcolorbox}
\section[merā nāṁ aman hai. merā kamm naukarī …]{(reading) — the same six facts, as sentences}

\label{lesson:PA-C45-pehla-paath}



\begin{quote}
Six labels, then six lines of a form. Now the same person, in sentences.

Forty-three words, and not one of them is new.
\end{quote}

\begin{tcolorbox}[breakable,skin=enhanced,colback=orange!6,colframe=orange!45!black,boxrule=0.5pt,arc=1mm,left=6pt,right=6pt,top=4pt,bottom=4pt,fonttitle=\bfseries,title={Read this}]
\begin{quote}
\pa{ਮੇਰਾ} \pa{ਨਾਂ} \pa{ਅਮਨ} \pa{ਹੈ।} \pa{ਮੇਰਾ} \pa{ਕੰਮ} \pa{ਨੌਕਰੀ} \pa{ਹੈ।} \pa{ਮੈਂ} \pa{ਪੰਜਾਬੀ} \pa{ਬੋਲਦਾ} \pa{ਹਾਂ।} \pa{ਮੈਂ} \pa{ਸੋਚਦਾ} \pa{ਹਾਂ} \pa{ਕਿ} \pa{ਪੰਜਾਬੀ} \pa{ਚੰਗਾ} \pa{ਹੈ।} \pa{ਮਨਨ} \pa{ਮੇਰਾ} \pa{ਦੋਸਤ} \pa{ਹੈ।} \pa{ਮਨਨ} \pa{ਦਾ} \pa{ਕੰਮ} \pa{ਖੇਤੀ} \pa{ਹੈ।} \pa{ਮਨਨ} \pa{ਠੀਕ}-\pa{ਠਾਕ} \pa{ਹੈ।} \pa{ਪਰ} \pa{ਮੈਨੂੰ} \pa{ਨਹੀਂ} \pa{ਪਤਾ} \pa{ਕਿ} \pa{ਉਹ} \pa{ਕਿੱਥੇ} \pa{ਹੈ।} \pa{ਅਸੀਂ} \pa{ਜਲਦੀ} \pa{ਮਿਲਾਂਗੇ।} \pa{ਧੰਨਵਾਦ।}
\end{quote}

Read it once without stopping. Do not translate as you go — let the sentences arrive.

Now read it again, and notice how little work it took.
\end{tcolorbox}

\begin{grammarlens}[title={What you've built}]
The form said the same things in six lines with no verb anywhere. Here \textbf{\pa{ਹੈ}} has to be said out loud, and it comes last every time.

\textbf{\pa{ਕਿ}} hangs a whole sentence off another one. \textbf{\pa{ਪਰ}} turns the paragraph. And after \textbf{\pa{ਮਨਨ}} is named twice, he becomes \textbf{\pa{ਉਹ}} — you followed that without being asked to.

This is the first whole paragraph the book has asked you to hold at once, and it works because every piece was paid for.
\end{grammarlens}

\subsection*{Your turn}
\begin{itemize}
\raggedright
  \item \emph{Say it:} the paragraph aloud, once, without stopping
\end{itemize}

\emph{Twice through:}

\begin{tcolorbox}[breakable,colback=teal!4,colframe=teal!35!black,title={Before you move on}]
How many words in that paragraph were new? (\textbf{None.})

What does the form leave out that the paragraph has to say? (\textbf{\pa{ਹੈ}.})
\end{tcolorbox}
